% Домашна работа 2 – работна програма на проекта YTKI
% Компилира се с pdfLaTeX (напр. в Overleaf) или с XeLaTeX.
\documentclass[11pt,a4paper]{article}

\usepackage{iftex}
\ifPDFTeX
  \usepackage[T2A]{fontenc}
  \usepackage[utf8]{inputenc}
  \usepackage[bulgarian]{babel}
\else
  \usepackage{fontspec}
  \usepackage{polyglossia}
  \setmainlanguage{bulgarian}
  \setmainfont{CMU Serif}
  \setmonofont{CMU Typewriter Text}
\fi

\usepackage[margin=2.2cm]{geometry}
\usepackage{array,tabularx,xltabular,booktabs,multirow}
\usepackage[table]{xcolor}
\usepackage{tikz}
\usetikzlibrary{arrows.meta,fit,shapes.geometric}
\usepackage{caption}
\usepackage{needspace}
\usepackage{enumitem}
\usepackage{pdflscape}
\usepackage{afterpage}
\usepackage{hyperref}
\hypersetup{colorlinks=true,linkcolor=black,urlcolor=black}

\setlength{\parindent}{0pt}
\setlength{\parskip}{5pt}
\setlength{\emergencystretch}{2em}
\renewcommand{\arraystretch}{1.2}
\definecolor{hdr}{gray}{0.90}
\definecolor{area}{gray}{0.96}
\setlist{nosep,topsep=2pt,leftmargin=*}

\captionsetup{labelfont=bf,labelsep=period,font=small,skip=6pt,hypcap=false}
\captionsetup[table]{position=top,justification=raggedright,singlelinecheck=false}

\newcolumntype{L}{>{\raggedright\arraybackslash}X}
\newcolumntype{P}[1]{>{\raggedright\arraybackslash}p{#1}}

% Ред със задача или резултат: номерът е отляво, а пренесеният текст
% се подравнява под началото на текста, не под номера.
% Всеки ред започва нов абзац; последният се затваря от клетката.
\newcommand{\tk}[2]{\par\hangindent=2.1em\hangafter=1\noindent
  \makebox[2.1em][l]{#1}#2}
\newcommand{\tkc}[2]{\par\hangindent=1.7em\hangafter=1\noindent
  \makebox[1.7em][l]{#1}#2}
\newcommand{\dl}[2]{\par\hangindent=2.9em\hangafter=1\noindent
  \makebox[2.9em][l]{#1}#2}

% Таблица 2 – описание на един работен пакет (по шаблона).
% #1 номер и име, #2 начало, #3 край, #4 цел, #5 задачи, #6 резултати
\newcommand{\wptable}[6]{%
  \par\noindent
  \begin{tabular}{|p{\dimexpr(\textwidth-4\tabcolsep-3\arrayrulewidth)/2\relax}
                  |p{\dimexpr(\textwidth-4\tabcolsep-3\arrayrulewidth)/2\relax}|}
    \hline
    \rowcolor{hdr}
    \multicolumn{2}{|>{\raggedright\arraybackslash}p{\dimexpr\textwidth-2\tabcolsep-2\arrayrulewidth\relax}|}%
      {\textbf{Номер и име на работен пакет:} #1}\\ \hline
    \textbf{Начало} (месец): #2 & \textbf{Край} (месец): #3\\ \hline
    \multicolumn{2}{|>{\raggedright\arraybackslash}p{\dimexpr\textwidth-2\tabcolsep-2\arrayrulewidth\relax}|}%
      {\textbf{Цел:} #4}\\ \hline
    \multicolumn{2}{|>{\raggedright\arraybackslash}p{\dimexpr\textwidth-2\tabcolsep-2\arrayrulewidth\relax}|}%
      {\textbf{Основни задачи:}#5}\\ \hline
    \multicolumn{2}{|>{\raggedright\arraybackslash}p{\dimexpr\textwidth-2\tabcolsep-2\arrayrulewidth\relax}|}%
      {\textbf{Резултати (deliverables):}#6}\\ \hline
  \end{tabular}\par\bigskip}

% ---------------------------------------------------------------------
% Стилове за диаграмите
\tikzset{
  % WBS
  wbsroot/.style={draw, fill=hdr, thick, inner sep=5pt, align=center,
    font=\small\bfseries},
  wbsarea/.style={draw, fill=area, rounded corners=2pt, inner sep=4pt},
  wphead/.style={draw, fill=hdr, inner sep=4pt, font=\footnotesize\bfseries},
  wptasks/.style={draw, inner sep=4pt, font=\scriptsize},
  % архитектура
  svc/.style={draw, rounded corners=2pt, fill=white, align=center,
    font=\scriptsize, minimum height=0.95cm, inner sep=2pt,
    text width=2.3cm},
  ext/.style={svc, fill=hdr},
  stp/.style={svc, text width=1.95cm},
  arr/.style={-{Stealth[length=4pt]}, line width=0.5pt},
  lbl/.style={font=\scriptsize\itshape, inner sep=1pt},
}
\newlength{\wpw}\setlength{\wpw}{2.9cm}
\newlength{\cardh}\setlength{\cardh}{3.85cm}
% Карта на работен пакет: #1 име на възел, #2 позиция, #3 заглавие, #4 задачи
\newcommand{\wpcard}[4]{%
  \node[wphead, anchor=north] (#1) at (#2)
    {\parbox[c][1.3cm][c]{\dimexpr\wpw-8pt\relax}{\centering #3}};
  \node[wptasks, anchor=north] (#1t) at (#1.south)
    {\parbox[t][\cardh][t]{\dimexpr\wpw-8pt\relax}{\raggedright
      \let\tk\tkc #4\par}};}
% Кутия за направление (ниво 1 в WBS): #1 име, #2 център, #3 ширина, #4 текст
\newcommand{\wbsarea}[4]{%
  \node[wbsarea, anchor=north] (#1) at (#2)
    {\parbox[c][1.6cm][c]{\dimexpr#3-8pt\relax}{\centering\footnotesize #4}};}

\begin{document}

% ---------------------------------------------------------------------
\begin{center}
  {\LARGE\bfseries YTKI: Интелигентна мултиагентна система\\[2pt]
  за мониторинг и реакция при инциденти в облачни услуги}\\[8pt]
  {\large Домашна работа 2: работна програма на проекта}\\[10pt]
  \begin{tabular}{ll}
    \toprule
    \textbf{Име и фамилия} & \textbf{Факултетен номер}\\ \midrule
    Йоан Байчев    & 0MI0600328\\
    Тугай Кафа     & 4MI0600275\\
    Калоян Марков  & 4MI0600392\\
    Илиан Запрянов & 8MI0600405\\ \bottomrule
  \end{tabular}
\end{center}

% Номерацията продължава след т. 1 от шаблона (Домашна работа 1).
\setcounter{section}{1}
\section{Описание на работната програма по проекта}

\subsection{Общо описание}

\textbf{Основна цел.} Да изградим и пуснем в употреба система от AI агенти,
която денонощно следи алъртите, логовете и health сигналите от всички среди,
сама прави първоначалното разследване на всеки инцидент и подава на дежурния
инженер готов анализ. Така искаме да съкратим времето за откриване
и възстановяване (MTTD и MTTR) и да намалим броя на фалшивите алърти, които
стигат до хората.

\textbf{Краен резултат.} Работеща в продукция версия 1.0 на YTKI, внедрена
в AWS (регион Франкфурт). Тя приема алъртите от dev, staging, int-prod и prod,
разследва ги с помощта на агенти, класифицира ги като фалшив алърт,
деградация или прекъсване, създава събитие в Splunk и отваря тикет към
отговорния екип. Към нея има табло, чат „Ask Alerts“ и документация за
потребителите и за поддръжката.

\textbf{Очакван ефект.} По-кратки и по-редки прекъсвания за клиентите,
по-малко излишни събуждания на дежурните и проследимост на инцидентите,
която улеснява докладването по NIS2.

\textbf{Подцели и междинни резултати.} За да стигнем до крайния резултат,
разделяме работата на шест подцели. Всяка от тях е цел на отделен работен
пакет и завършва с междинен резултат (milestone), който може да се провери.
Към тях се добавят задължителните пакети за комуникации (РП7) и за управление
на проекта (РП8).

\begin{center}\small
\begin{tabularx}{\textwidth}{|c|L|L|c|}
  \hline
  \rowcolor{hdr}
  \textbf{РП} & \textbf{Подцел} & \textbf{Междинен резултат} & \textbf{Месец}\\ \hline
  1 & Да изясним нуждите, изискванията и ограниченията и да проектираме
      системата и начина, по който ще я внедрим в AWS
    & \textbf{M1.} Изисквания и архитектура, одобрени от DevOps и SRE екипите
    & 2\\ \hline
  2 & Да изградим средата в AWS и алъртите от всички среди да постъпват
      в реално време и в общ формат
    & \textbf{M2.} Алъртите от четирите среди постъпват в YTKI, разгърната в AWS
    & 5\\ \hline
  3 & Агентите сами да разследват всеки нов алърт
    & \textbf{M3.} За всеки алърт се създава отчет с хипотеза за първопричина
    & 8\\ \hline
  4 & Класификацията да е достатъчно точна, за да могат дежурните да разчитат
      на нея
    & \textbf{M4.} F1 поне 0,8 върху тестов набор от минали инциденти
    & 9\\ \hline
  5 & Резултатите от анализа да стигат автоматично до отговорния екип
    & \textbf{M5.} Алърт минава целия път до събитие в Splunk и тикет
    & 10\\ \hline
  \multirow{2}{*}{6}
    & \multirow{2}{=}{Системата да бъде проверена в реални условия
      и пусната в продукция}
    & \textbf{M6.} Приключен пилот в сенчест режим в staging и int-prod & 11\\ \cline{3-4}
  & & \textbf{M7.} YTKI v1.0 работи в prod & 12\\ \hline
\end{tabularx}
\end{center}

\textbf{Продължителност и етапи.} Проектът е планиран за 12 месеца. Първите
два месеца са за иницииране и планиране. Основната разработка е от втори
до десети месец, а от осми до единадесети тестваме и провеждаме пилот.
Последният месец е за пускане в продукция и приключване на проекта.
Управлението и комуникациите вървят през цялото време. Разпределението
на пакетите във времето и междинните резултати са показани на
фиг.~\ref{fig:gantt}.

\begin{center}
\begin{tikzpicture}[x=0.9cm,y=0.62cm]
  \foreach \m in {1,...,12}
    \node[font=\scriptsize] at (\m-0.5,0.45) {\m};
  \node[font=\scriptsize\itshape, anchor=west] at (-5.4,0.45) {месец};
  \foreach \m in {0,...,12}
    \draw[gray!35] (\m,0.1) -- (\m,-8.5);
  \foreach \i/\s/\e/\c/\name in {
      1/1/2/45/{РП1. Анализ и проектиране},
      2/2/5/45/{РП2. Облачна платформа и приемане},
      3/3/8/45/{РП3. Агентно разследване},
      4/4/10/45/{РП4. Класификация и самообучение},
      5/6/10/45/{РП5. Интеграции и интерфейс},
      6/8/12/45/{РП6. Тестване и внедряване},
      7/1/12/22/{РП7. Комуникации},
      8/1/12/22/{РП8. Управление на проекта}} {
    \node[font=\scriptsize, anchor=west] at (-5.4,-\i) {\name};
    \fill[gray!\c] (\s-1,-\i-0.24) rectangle (\e,-\i+0.24);
    \draw[gray!70] (\s-1,-\i-0.24) rectangle (\e,-\i+0.24);
  }
  \foreach \i/\m/\name in {1/2/M1,2/5/M2,3/8/M3,4/9/M4,5/10/M5,6/11/M6,6/12/M7} {
    \node[diamond, fill=black, inner sep=1.5pt] at (\m,-\i) {};
    \node[font=\tiny\bfseries, anchor=south] at (\m,-\i+0.27) {\name};
  }
\end{tikzpicture}
\captionof{figure}{Разпределение на работните пакети във времето и междинни
резултати (M1--M7)}
\label{fig:gantt}
\end{center}

% ---------------------------------------------------------------------
\Needspace{8\baselineskip}
\subsection{Работна програма и задачи}

\subsubsection*{Йерархична структура на работата}

Структурата на работата е организирана по резултати (deliverable-based WBS)
и е показана на фиг.~\ref{fig:wbs}. На първо ниво работните пакети са
групирани в четири направления според експертизата, която изискват: анализ
и проектиране (бизнес анализатори и архитект), имплементиране и внедряване
(разработчици, QA и DevOps инженери), комуникации и управление на проекта.
Всяко направление завършва със собствен резултат, а проектът е готов, когато
са готови всички. На второ ниво са работните пакети, а на трето~-- задачите
в тях.

\subsubsection*{Обща методология}

Работата е разделена така, че всеки пакет да събира дейности, за които
е нужна сходна експертиза, и да завършва с резултат, който може да се провери.
Имплементирането е най-голямото направление, затова го разделяме на пет
пакета. Четири от тях следват пътя на един алърт през системата
(фиг.~1 от Домашна работа~1): приемане, разследване, класификация
и оперативно действие. Петият обхваща тестването и пускането в продукция.

В РП1 работят бизнес анализатор и архитект. В РП2 основната роля е на DevOps
и backend инженерите, в РП3 са нужни разработчици с опит с LLM и агентни
системи, а в РП4 работи специалист по данни. РП5 е за разработчиците на
интеграциите и интерфейса, а РП6 е за QA и SRE инженерите.

Пакетите от РП2 до РП5 се застъпват във времето. Работим на двуседмични
спринтове и всеки пакет може да започне, щом предишният е дал първа работеща
версия. Част от РП3 и РП5 стъпва на вече работещия прототип (таблото
и „Ask Alerts“), така че там голяма част от работата е разширяване,
а не разработка от нулата.

Начинът, по който системата ще бъде внедрена в AWS, е описан в
Приложение~А. Архитектурата се уточнява в задача 1.6, средата се изгражда
в РП2, а пускането в продукция е в РП6.

% WBS на отделна хоризонтална страница; текстът продължава преди нея.
\afterpage{%
\begin{landscape}
\centering
\vspace*{\fill}
\begin{tikzpicture}[line width=0.5pt]
  % центрове на осемте колони (ширина на колона 2.9 cm, разстояние 0.3 cm)
  \foreach \i in {1,...,8}
    \coordinate (c\i) at ({1.45 + (\i-1)*3.2}, 0);
  \coordinate (mid) at (12.65,0);
  \coordinate (impl) at (11.05,0);

  % ниво 0
  \node[wbsroot, anchor=north] (root) at (mid) {YTKI v1.0\\[1pt]
    \normalfont\footnotesize мултиагентна система за мониторинг и реакция
    при инциденти в облачни услуги};

  % ниво 1: направления
  \coordinate (a) at ([yshift=-0.9cm]root.south);
  \wbsarea{ar1}{c1 |- a}{\wpw}{\textbf{Анализ и проектиране}\\
    {\scriptsize\itshape анализатори, архитект}}
  \wbsarea{ar2}{impl |- a}{15.7cm}{\textbf{Имплементиране и внедряване}\\
    \textit{разработчици, QA, DevOps и SRE инженери}}
  \wbsarea{ar3}{c7 |- a}{\wpw}{\textbf{Комуникации}\\
    {\scriptsize\itshape комуникации и маркетинг}}
  \wbsarea{ar4}{c8 |- a}{\wpw}{\textbf{Управление}\\
    {\scriptsize\itshape ръководител на проекта}}

  % ниво 2 и 3: работни пакети и задачи
  \coordinate (w) at ([yshift=-0.8cm]ar2.south);
  \wpcard{wp1}{c1 |- w}{РП1. Анализ и проектиране}{%
    \tk{1.1}{Интервюта}
    \tk{1.2}{Анкетно проучване}
    \tk{1.3}{Анализ на алъртите}
    \tk{1.4}{Изисквания и метрики}
    \tk{1.5}{Данни и съответствие}
    \tk{1.6}{Архитектура в AWS}}
  \wpcard{wp2}{c2 |- w}{РП2. Облачна платформа и приемане}{%
    \tk{2.1}{AWS акаунти и мрежа}
    \tk{2.2}{CI/CD конвейер}
    \tk{2.3}{Входна точка за алъртите}
    \tk{2.4}{Модел на събитието и разпращане}
    \tk{2.5}{Свързване на средите}
    \tk{2.6}{Маскиране на данни}}
  \wpcard{wp3}{c3 |- w}{РП3. Агентно разследване}{%
    \tk{3.1}{Оркестрация и планиращ агент}
    \tk{3.2}{Агент за логове}
    \tk{3.3}{Агент за health проверки}
    \tk{3.4}{Агент за зависимостите}
    \tk{3.5}{Хипотеза за първопричина}}
  \wpcard{wp4}{c4 |- w}{РП4. Класификация и самообучение}{%
    \tk{4.1}{Маркиране на инциденти}
    \tk{4.2}{Класификатор}
    \tk{4.3}{Playwright сценарии}
    \tk{4.4}{Класически модели}
    \tk{4.5}{Обратна връзка}}
  \wpcard{wp5}{c5 |- w}{РП5. Интеграции и интерфейс}{%
    \tk{5.1}{Мапинг на компонентите}
    \tk{5.2}{Събития в Splunk}
    \tk{5.3}{Тикети}
    \tk{5.4}{Потоци в n8n}
    \tk{5.5}{Табло и „Ask Alerts“}}
  \wpcard{wp6}{c6 |- w}{РП6. Тестване и внедряване}{%
    \tk{6.1}{Тестове на целия поток}
    \tk{6.2}{Натоварване}
    \tk{6.3}{Сигурност}
    \tk{6.4}{Пилот в сенчест режим}
    \tk{6.5}{Пускане в prod}
    \tk{6.6}{Наблюдение на YTKI}}
  \wpcard{wp7}{c7 |- w}{РП7. Комуникации}{%
    \tk{7.1}{Комуникационен план}
    \tk{7.2}{Информационни материали}
    \tk{7.3}{Демонстрации}
    \tk{7.4}{Обучителни уъркшопи}
    \tk{7.5}{Финална презентация}}
  \wpcard{wp8}{c8 |- w}{РП8. Управление на проекта}{%
    \tk{8.1}{План на проекта}
    \tk{8.2}{Следене на напредъка}
    \tk{8.3}{Финанси}
    \tk{8.4}{Човешки ресурси}
    \tk{8.5}{Заинтересовани лица}
    \tk{8.6}{Рискове}
    \tk{8.7}{Приключване}}

  % връзки
  \coordinate (b0) at ([yshift=-0.45cm]root.south);
  \draw (root.south) -- (b0);
  \draw (c1 |- b0) -- (c8 |- b0);
  \foreach \n/\x in {ar1/c1, ar2/impl, ar3/c7, ar4/c8}
    \draw (\x |- b0) -- (\n.north);
  \coordinate (b1) at ([yshift=-0.4cm]ar2.south);
  \draw (ar2.south) -- (b1);
  \draw (c2 |- b1) -- (c6 |- b1);
  \foreach \i in {2,...,6} \draw (c\i |- b1) -- (wp\i.north);
  \foreach \i/\a in {1/ar1, 7/ar3, 8/ar4} \draw (\a.south) -- (wp\i.north);
\end{tikzpicture}
\captionof{figure}{Йерархична структура на работата по проекта
(deliverable-based WBS): направления, работни пакети и задачи}
\label{fig:wbs}
\vspace*{\fill}
\end{landscape}}

\Needspace{24\baselineskip}
\subsubsection*{Обща структура на работната програма}

\begin{xltabular}{\textwidth}{|c|P{4cm}|L|}
  \caption{Работни пакети}\label{tab:wp}\\
  \hline
  \rowcolor{hdr}
  \textbf{Номер} & \textbf{Име на работен пакет} & \textbf{Цел на работен пакет}\\ \hline
  \endfirsthead
  \hline
  \rowcolor{hdr}
  \textbf{Номер} & \textbf{Име на работен пакет} & \textbf{Цел на работен пакет}\\ \hline
  \endhead
  \multicolumn{3}{|l|}{\textit{Анализ и проектиране}}\\ \hline
  РП1 & Анализ и проектиране
      & (подцел 1) Изясняване на нуждите на екипите, изискванията, наличните
        данни и ограниченията и проектиране на системата и на внедряването
        в AWS\\ \hline
  \multicolumn{3}{|l|}{\textit{Имплементиране и внедряване}}\\ \hline
  РП2 & Облачна платформа и приемане на данни
      & (подцел 2) Изграждане на средата в AWS и единен поток в реално време
        на алъртите, логовете и health сигналите от всички наблюдавани среди\\ \hline
  РП3 & Агентно разследване
      & (подцел 3) Автоматично разследване на всеки нов алърт от
        специализирани агенти\\ \hline
  РП4 & Класификация и самообучение
      & (подцел 4) Точна класификация на инцидентите, която се подобрява
        с натрупването на история\\ \hline
  РП5 & Интеграции и интерфейс
      & (подцел 5) Резултатите от анализа да стигат автоматично до
        отговорния екип и да се виждат на едно място\\ \hline
  РП6 & Тестване и внедряване
      & (подцел 6) Проверка на системата в реални условия и пускане
        в продукция\\ \hline
  \multicolumn{3}{|l|}{\textit{Комуникации}}\\ \hline
  РП7 & Комуникации
      & Заинтересованите лица да знаят за проекта, да разбират резултатите
        и да могат да работят със системата\\ \hline
  \multicolumn{3}{|l|}{\textit{Управление}}\\ \hline
  РП8 & Управление на проекта
      & Проектът да приключи в срок, в рамките на бюджета и с нужното
        качество\\ \hline
\end{xltabular}

\Needspace{30\baselineskip}
\subsubsection*{Описание на работните пакети}

{\captionof{table}{Описание на работните пакети}}

\wptable{РП1: Анализ и проектиране}{1}{2}
{Изясняване на нуждите на екипите, изискванията, наличните данни
и ограниченията и проектиране на системата и на внедряването в AWS.}
{\tk{1.1.}{Интервюта с дежурните инженери и собствениците на услуги}
 \tk{1.2.}{Анкетно проучване сред DevOps и SRE екипите (брой алърти,
   фалшиви тревоги, време за реакция)}
 \tk{1.3.}{Анализ на алъртите и инцидентите, събрани от прототипа}
 \tk{1.4.}{Изисквания към системата и метрики за успех (MTTD, MTTR, дял
   на фалшивите алърти)}
 \tk{1.5.}{Правила за работа с данните и съответствие (GDPR, NIS2, какви
   данни може да се подават към LLM)}
 \tk{1.6.}{Архитектура на системата и план за внедряване в AWS
   (Приложение~А)}}
{\dl{D1.1.}{Обобщение на интервютата (м.~1)}
 \dl{D1.2.}{Доклад от анкетата (м.~1)}
 \dl{D1.3.}{Доклад за текущия процес с базови стойности на MTTD и MTTR (м.~2)}
 \dl{D1.4.}{Спецификация на изискванията (м.~2)}
 \dl{D1.5.}{Правила за работа с данните (м.~2)}
 \dl{D1.6.}{Архитектурен документ и план за внедряване (м.~2)}}

\wptable{РП2: Облачна платформа и приемане на данни}{2}{5}
{Изграждане на средата в AWS и единен поток в реално време на алъртите,
логовете и health сигналите от всички наблюдавани среди.}
{\tk{2.1.}{AWS акаунти, мрежа и базова сигурност, описани като код
   с Terraform}
 \tk{2.2.}{CI/CD конвейер за всички компоненти}
 \tk{2.3.}{Входна точка за алъртите (API Gateway, удостоверяване на
   източниците, парсери за основните формати)}
 \tk{2.4.}{Общ модел на събитието и разпращане към подписчиците чрез
   EventBridge}
 \tk{2.5.}{Свързване на четирите наблюдавани среди, включително prod}
 \tk{2.6.}{Маскиране на личните данни, преди да се подадат към LLM}}
{\dl{D2.1.}{Dev и staging среди на YTKI в AWS (м.~3)}
 \dl{D2.2.}{Работещ CI/CD конвейер (м.~3)}
 \dl{D2.3.}{Услуга за приемане на алърти (м.~4)}
 \dl{D2.4.}{Модел на събитието и шина за разпращане (м.~4)}
 \dl{D2.5.}{Работещи връзки към четирите среди (м.~5)}
 \dl{D2.6.}{Модул за маскиране на данни (м.~5)}}

\wptable{РП3: Агентно разследване}{3}{8}
{Автоматично разследване на всеки нов алърт от специализирани агенти.}
{\tk{3.1.}{Работен поток на разследването в Step Functions и планиращ агент,
   който определя какво да се провери и разпределя подзадачите}
 \tk{3.2.}{Агент за търсене и анализ на логове в Cloud Logging}
 \tk{3.3.}{Агент за проверка на health endpoint-ите}
 \tk{3.4.}{Агент за графа на зависимостите и оценка на каскадния риск}
 \tk{3.5.}{Обединяване на резултатите от агентите в хипотеза за първопричина}}
{\dl{D3.1.}{Работен поток с планиращ агент (м.~5)}
 \dl{D3.2.}{Агент за логове (м.~6)}
 \dl{D3.3.}{Агент за health проверки (м.~6)}
 \dl{D3.4.}{Агент за зависимостите (м.~7)}
 \dl{D3.5.}{Отчет за разследване, създаван за всеки алърт (м.~8)}}

\wptable{РП4: Класификация и самообучение}{4}{10}
{Точна класификация на инцидентите, която се подобрява с натрупването на
история.}
{\tk{4.1.}{Събиране и ръчно маркиране на минали инциденти}
 \tk{4.2.}{Класификация по трите категории с вероятност за всяка}
 \tk{4.3.}{Playwright сценарии за потвърждаване на деградация или прекъсване}
 \tk{4.4.}{Обучение и сравнение на логистична регресия и random forest}
 \tk{4.5.}{Обратна връзка от дежурните и периодично преобучение}}
{\dl{D4.1.}{Маркиран набор от исторически инциденти (м.~5)}
 \dl{D4.2.}{Първа версия на класификатора (м.~7)}
 \dl{D4.3.}{Playwright сценарии за ключовите услуги (м.~8)}
 \dl{D4.4.}{Обучени модели и доклад от сравнението (м.~9)}
 \dl{D4.5.}{Работещ цикъл за самообучение (м.~10)}}

\wptable{РП5: Интеграции и интерфейс}{6}{10}
{Резултатите от анализа да стигат автоматично до отговорния екип и да се
виждат на едно място.}
{\tk{5.1.}{Мапинг между алърт, компонент и отговорен екип}
 \tk{5.2.}{Автоматично създаване на събития в Splunk}
 \tk{5.3.}{Автоматично отваряне на тикети с приложен анализ}
 \tk{5.4.}{Потоци в n8n за най-честите типове инциденти}
 \tk{5.5.}{Добавяне на класификацията и хипотезите в таблото и „Ask Alerts“}}
{\dl{D5.1.}{Таблица за мапинг на компонентите (м.~7)}
 \dl{D5.2.}{Интеграция със Splunk (м.~8)}
 \dl{D5.3.}{Интеграция с тикет системата (м.~9)}
 \dl{D5.4.}{Набор от n8n потоци (м.~10)}
 \dl{D5.5.}{Обновени табло и „Ask Alerts“ (м.~10)}}

\wptable{РП6: Тестване и внедряване}{8}{12}
{Проверка на системата в реални условия и пускане в продукция.}
{\tk{6.1.}{Тестване на целия поток с исторически и симулирани алърти}
 \tk{6.2.}{Тестове за натоварване (алърт бури) и отказоустойчивост}
 \tk{6.3.}{Преглед на сигурността и защитата на данните}
 \tk{6.4.}{Пилот в сенчест режим в staging и int-prod: YTKI анализира
   алъртите, но не отваря тикети}
 \tk{6.5.}{Изграждане на prod средата в AWS и поетапно пускане}
 \tk{6.6.}{Наблюдение на самата YTKI: аларми, табла и контролен сигнал}}
{\dl{D6.1.}{Отчет от функционалните тестове (м.~9)}
 \dl{D6.2.}{Отчет от тестовете за натоварване (м.~10)}
 \dl{D6.3.}{Доклад от прегледа на сигурността (м.~10)}
 \dl{D6.4.}{Пилотен отчет с измерените метрики (м.~11)}
 \dl{D6.5.}{YTKI v1.0 в продукция (м.~12)}
 \dl{D6.6.}{Табла, аларми и runbook-ове за поддръжката (м.~12)}}

\wptable{РП7: Комуникации}{1}{12}
{Заинтересованите лица да знаят за проекта, да разбират резултатите и да
могат да работят със системата.}
{\tk{7.1.}{Изготвяне на комуникационен план}
 \tk{7.2.}{Информационни материали: кратко ръководство за потребителя,
   видео демонстрация, страница във вътрешното уики}
 \tk{7.3.}{Демонстрации пред екипите след всеки междинен резултат}
 \tk{7.4.}{Обучителни уъркшопи за дежурните екипи преди пускането в prod}
 \tk{7.5.}{Финална презентация пред ръководството}}
{\dl{D7.1.}{Комуникационен план (м.~1)}
 \dl{D7.2.}{Ръководство за потребителя, видео и страница в уикито (м.~10)}
 \dl{D7.3.}{Проведени демонстрации с кратко обобщение след всяка (м.~2--12)}
 \dl{D7.4.}{Проведени уъркшопи и материали от тях (м.~11)}
 \dl{D7.5.}{Финална презентация (м.~12)}}

\wptable{РП8: Управление на проекта}{1}{12}
{Проектът да приключи в срок, в рамките на бюджета и с нужното качество.}
{\tk{8.1.}{Изготвяне на плана на проекта (обхват, график, бюджет)}
 \tk{8.2.}{Следене на напредъка чрез седмични срещи и месечни отчети}
 \tk{8.3.}{Управление на финансите, включително разходите за AWS и LLM}
 \tk{8.4.}{Управление на човешките ресурси: роли, натоварване, заместване
   при отсъствие}
 \tk{8.5.}{Управление на заинтересованите лица}
 \tk{8.6.}{Управление на рисковете и промените}
 \tk{8.7.}{Приключване на проекта}}
{\dl{D8.1.}{План за управление на проекта (м.~1)}
 \dl{D8.2.}{Месечни отчети за напредъка (м.~1--12)}
 \dl{D8.3.}{Бюджет и месечен отчет за разходите (м.~1--12)}
 \dl{D8.4.}{Матрица на отговорностите (RACI) (м.~1)}
 \dl{D8.5.}{Регистър на заинтересованите лица и план за работа с тях (м.~1)}
 \dl{D8.6.}{Регистър на рисковете (от м.~1, поддържа се до края)}
 \dl{D8.7.}{Финален отчет с изводи за бъдещи проекти (м.~12)}}

% =====================================================================
\clearpage
\section*{Приложение А. Архитектура на внедряване в AWS}
\addcontentsline{toc}{section}{Приложение А. Архитектура на внедряване в AWS}

Това приложение описва как YTKI ще работи в AWS. То е основа за задача 1.6
(архитектура и план за внедряване), за РП2 (изграждане на средата) и за
задача 6.5 (пускане в продукция).

\subsection*{А.1. Какво изисква системата от инфраструктурата}

Изборът на архитектура следва от пет особености на YTKI.

\begin{itemize}
  \item \textbf{Неравномерно натоварване.} В спокойни часове идват по няколко
    алърта, а при голям инцидент могат да дойдат стотици за минути (т.нар.
    алърт буря). Системата трябва да поема пиковете, без да плащаме за
    капацитет, който стои неизползван.
  \item \textbf{Дълги задачи.} Разследването на един инцидент отнема от
    няколко секунди до няколко минути, защото агентите чакат LLM и външни
    системи. Това не бива да се прави синхронно в рамките на една HTTP заявка.
  \item \textbf{Надеждност.} YTKI работи 24/7 и трябва да е поне толкова
    надеждна, колкото системите, които наблюдава. Ако тя спре, алъртите
    не бива да се губят.
  \item \textbf{Чувствителни данни.} Логовете може да съдържат лични данни,
    затова данните остават в ЕС и се маскират, преди да стигнат до LLM.
  \item \textbf{Малък екип.} Нямаме капацитет да поддържаме сървъри или
    Kubernetes клъстер.
\end{itemize}

\subsection*{А.2. Избор на подход}

Разгледахме четири подхода (табл.~\ref{tab:approach}).

\begingroup\small
\begin{xltabular}{\textwidth}{|P{2.9cm}|L|L|}
  \caption{Сравнение на подходите за внедряване}\label{tab:approach}\\
  \hline
  \rowcolor{hdr}
  \textbf{Подход} & \textbf{Предимства} & \textbf{Недостатъци за YTKI}\\ \hline
  \endfirsthead
  \hline
  \rowcolor{hdr}
  \textbf{Подход} & \textbf{Предимства} & \textbf{Недостатъци за YTKI}\\ \hline
  \endhead
  Serverless (Lambda, Step Functions)
    & Плащаме само за реална работа; мащабира се автоматично при алърт бури;
      няма сървъри за поддръжка; висока наличност по подразбиране
    & До 15 минути на едно изпълнение на функция; студени стартове; по-трудно
      локално тестване\\ \hline
  Контейнери в ECS Fargate
    & Няма ограничение във времето; подходящо за браузъри и постоянно работещи
      услуги
    & Плащаме за задачите, докато работят, дори без натоварване; мащабирането
      се настройва ръчно\\ \hline
  Kubernetes (EKS)
    & Голяма гъвкавост и преносимост между облаци
    & Висока сложност и постоянна цена на клъстера; нужен е човек само за
      поддръжката му\\ \hline
  Виртуални машини (EC2)
    & Пълен контрол върху средата
    & Обновяване на ОС, ръчна висока наличност и мащабиране\\ \hline
\end{xltabular}
\endgroup

Избираме \textbf{serverless като основа} и ECS Fargate само за двата
компонента, които не пасват на Lambda: Playwright сценариите (нужен е цял
браузър) и n8n (постоянно работеща услуга). Разследването не се изпълнява
синхронно зад API-то. Входната точка само приема алърта, връща отговор
веднага и го оставя в опашка, а самото разследване се оркестрира от
Step Functions. Така времето за отговор към системите за мониторинг не
зависи от това колко бавно отговаря LLM.

\subsection*{А.3. Обща схема}

Схемата е показана на фиг.~\ref{fig:aws}. Системата има три входа: алъртите
от наблюдаваните среди (горе), потребителите на таблото и чата (долу)
и резервния път, по който алъртите стигат до дежурните и без YTKI.

% Архитектурна схема на отделна хоризонтална страница
\afterpage{%
\begin{landscape}
\centering
\vspace*{\fill}
\begin{tikzpicture}[line width=0.5pt]
  % ---------------- път на алъртите (горе)
  \node[ext] (src) at (1.5,12.2) {\textbf{Наблюдавани среди}\\dev, staging,
    int-prod, prod};
  \node[svc] (apiin) at (4.4,12.2) {\textbf{API Gateway}\\приемане на алърти};
  \node[svc] (lin) at (7.2,12.2) {\textbf{Lambda}\\парсване, дедупликация};
  \node[svc] (sqs) at (10.0,12.2) {\textbf{SQS}\\буфер и DLQ};
  \node[svc] (eb) at (12.8,12.2) {\textbf{EventBridge}\\разпращане};
  \node[ext] (pd) at (23.2,12.2) {\textbf{PagerDuty / Opsgenie}\\дежурни};
  \draw[arr] (src) -- (apiin);
  \draw[arr] (apiin) -- (lin);
  \draw[arr] (lin) -- (sqs);
  \draw[arr] (sqs) -- (eb);
  \draw[arr, dashed] (src.north) -- ++(0,0.65) -| (pd.north)
    node[pos=0.25, above, lbl] {резервен път: алъртите стигат до дежурните
    и без YTKI};

  % ---------------- VPC
  \draw[dashed, rounded corners=4pt] (6.2,3.85) rectangle (21.0,11.05);
  \node[lbl, anchor=north east] at (20.9,11.0)
    {VPC: частни подмрежи в 3 зони на достъпност};

  % ---------------- Step Functions
  \draw[rounded corners=3pt, line width=0.8pt] (6.6,4.2) rectangle (17.5,10.45);
  \node[font=\scriptsize, anchor=north west] at (6.7,10.4)
    {\textbf{Step Functions}: разследване на един инцидент};
  \node[stp] (plan) at (8.0,7.4) {\textbf{Планиране}\\Lambda + LLM};
  \node[stp] (alog) at (10.6,9.1) {\textbf{Логове}\\Lambda};
  \node[stp] (ahl)  at (10.6,7.95) {\textbf{Health}\\Lambda};
  \node[stp] (adep) at (10.6,6.8) {\textbf{Зависимости}\\Lambda};
  \node[stp] (apw)  at (10.6,5.65) {\textbf{Playwright}\\ECS Fargate};
  \node[draw, dotted, fit=(alog)(apw), inner sep=3pt] (par) {};
  \node[lbl, anchor=north, yshift=-2pt] at (par.south) {паралелно};
  \node[stp] (agg) at (13.2,7.4) {\textbf{Обединяване}\\хипотеза за причина};
  \node[stp] (cls) at (15.75,7.4) {\textbf{Класификация}\\модел + LLM};
  \node[stp] (act) at (15.75,5.3) {\textbf{Действие}\\Splunk, тикет};
  \foreach \a in {alog,ahl,adep,apw} {
    \draw[arr] (plan.east) -- (\a.west);
    \draw[arr] (\a.east) -- (agg.west);
  }
  \draw[arr] (agg) -- (cls);
  \draw[arr] (cls) -- (act);
  \draw[arr] (eb.south) -- (eb.south |- 0,10.45)
    node[pos=0.22, right, lbl] {старт на разследване};

  % ---------------- данни в VPC
  \node[svc] (nat) at (19.25,9.8) {\textbf{NAT Gateway}\\изход към интернет};
  \node[svc] (cache) at (19.25,8.45) {\textbf{ElastiCache (Valkey)}\\общ контекст};
  \node[svc] (aur) at (19.25,7.1) {\textbf{Aurora PostgreSQL}\\история, pgvector};
  \node[svc] (n8n) at (19.25,5.3) {\textbf{n8n}\\ECS Fargate};
  \draw[arr] (17.5,9.8) -- (nat.west);
  \draw[arr, <->] (17.5,8.45) -- (cache.west);
  \draw[arr, <->] (17.5,7.1) -- (aur.west);
  \draw[arr] (act.east) -- (n8n.west);

  % ---------------- външни системи
  \node[lbl] at (23.2,10.75) {Външни системи};
  \node[ext] (gcp) at (23.2,9.8) {\textbf{GCP Cloud Logging}\\логове на средите};
  \node[ext] (spl) at (23.2,8.45) {\textbf{Splunk}\\събития};
  \node[ext] (tix) at (23.2,7.1) {\textbf{Тикет система}};
  \draw (nat.east) -- (21.5,9.8);
  \draw[arr] (21.5,9.8) -- (gcp.west);
  \draw[arr] (21.5,9.8) |- (spl.west);
  \draw[arr] (21.5,9.8) |- (tix.west);

  % ---------------- регионални услуги
  \node[lbl, align=center] at (3.3,2.7) {Регионални услуги\\(от VPC през\\
    VPC endpoints)};
  \node[svc] (s3)  at (7.0,2.7) {\textbf{S3}\\табло и архив};
  \node[svc] (ddb) at (9.7,2.7) {\textbf{DynamoDB}\\инциденти, статус};
  \node[svc] (br)  at (12.4,2.7) {\textbf{Bedrock}\\LLM в ЕС регион};
  \node[svc] (sec) at (15.1,2.7) {\textbf{Secrets Manager}\\и KMS};
  \node[svc] (cw)  at (17.8,2.7) {\textbf{CloudWatch}\\X-Ray, аларми};
  \draw (7.0,3.45) -- (17.8,3.45);
  \foreach \n in {s3,ddb,br,sec,cw} \draw (\n.north) -- (\n.north |- 0,3.45);
  \draw[arr, <->] (11.05,4.2) -- (11.05,3.45);
  \node[lbl, anchor=west] at (11.15,3.62) {VPC endpoints};

  % ---------------- достъп на потребителите (долу)
  \node[ext] (usr) at (1.5,0.75) {\textbf{Потребители}\\DevOps, SRE, дежурни};
  \node[svc] (r53) at (4.3,0.75) {\textbf{Route 53}\\DNS};
  \node[svc] (cf)  at (7.0,0.75) {\textbf{CloudFront}\\и AWS WAF};
  \node[svc] (apiui) at (9.7,0.75) {\textbf{API Gateway}\\API и чат (стрийминг)};
  \node[svc] (lui) at (12.4,0.75) {\textbf{Lambda}\\API и „Ask Alerts“};
  \node[svc] (cog) at (15.1,0.75) {\textbf{Cognito}\\вход през SSO};
  \draw[arr] (usr) -- (r53);
  \draw[arr] (r53) -- (cf);
  \draw[arr] (cf) -- (apiui);
  \draw[arr] (cf.north) -- (s3.south);
  \draw[arr] (apiui) -- (lui);
  \draw[arr] (lui.north) -- (br.south);
  \draw[arr] (lui.north west) -- (ddb.south east);
  \draw[dashed] (lui) -- (cog);
\end{tikzpicture}
\captionof{figure}{Архитектура на YTKI в AWS. Сивите кутии са външни системи.
Step Functions е регионална услуга; в VPC работят функциите и задачите,
които тя стартира.}
\label{fig:aws}
\vspace*{\fill}
\end{landscape}}

\subsection*{А.4. Път на един алърт през системата}

\begin{enumerate}
  \item Системата за мониторинг (напр. Alertmanager или Grafana) изпраща
    алърта едновременно към PagerDuty или Opsgenie и към входната точка на
    YTKI. Така дежурните получават алърта дори ако YTKI не работи.
  \item API Gateway проверява токена на източника и ограничава броя на
    заявките. Lambda функцията парсва алърта, превръща го в общия модел на
    събитието и изчислява неговия отпечатък (fingerprint). Ако същият алърт
    вече е приет в последните минути, той се отбелязва като дубликат
    и не се разследва отново.
  \item Алъртът влиза в SQS опашка. Тя поема алърт бурите, а съобщенията,
    които не могат да се обработят, отиват в отделна опашка (DLQ) за
    преглед и повторно изпращане.
  \item EventBridge разпраща събитието към абонатите. Това е абонаментният
    модел от фиг.~1 на Домашна работа~1: правило стартира разследване,
    друго записва събитието в архива, а в бъдеще може да се добавят нови
    абонати без промяна на кода.
  \item Step Functions изпълнява разследването. Планиращият агент решава
    какво да се провери, след което агентите за логове, health проверки,
    зависимости и Playwright работят паралелно. Междинните им резултати се
    записват в общ контекст в ElastiCache.
  \item Резултатите се обединяват в хипотеза за първопричина. Класификаторът
    определя дали става дума за фалшив алърт, деградация или прекъсване,
    с вероятност за всяка категория. Историята на подобни минали инциденти
    се търси в Aurora (pgvector).
  \item При деградация или прекъсване системата създава събитие в Splunk,
    отваря тикет към отговорния екип и при нужда стартира поток в n8n.
    Отчетът за разследването се записва в DynamoDB и S3 и се появява
    в таблото.
  \item Когато дежурният потвърди или поправи класификацията, отговорът му
    се записва като нов пример за обучение (РП4).
\end{enumerate}

\subsection*{А.5. Компоненти}

\begingroup\small
\begin{xltabular}{\textwidth}{|P{2.6cm}|P{4.1cm}|L|}
  \caption{Компоненти на системата и роля на всеки от тях}\label{tab:aws}\\
  \hline
  \rowcolor{hdr}
  \textbf{Слой} & \textbf{AWS услуга} & \textbf{За какво я използваме}\\ \hline
  \endfirsthead
  \hline
  \rowcolor{hdr}
  \textbf{Слой} & \textbf{AWS услуга} & \textbf{За какво я използваме}\\ \hline
  \endhead
  DNS & Route 53
    & Домейните на таблото и на входната точка за алъртите\\ \hline
  Край на мрежата & CloudFront и AWS WAF
    & Доставят таблото, филтрират злонамерен трафик и ограничават броя
      на заявките\\ \hline
  Удостоверяване & Cognito
    & Вход през корпоративния доставчик на идентичност (SSO), JWT токени
      за API-то\\ \hline
  API & API Gateway (REST, регионален)
    & Две API-та: за приемане на алърти и за таблото и чата. Отговорите
      в чата се стриймват към потребителя\\ \hline
  Буфер & SQS с DLQ
    & Поема алърт бурите и пази съобщенията, които не са обработени\\ \hline
  Разпращане & EventBridge
    & Абонаментният модел: правила насочват събитията към абонатите;
      архив за повторно изпращане\\ \hline
  Оркестрация & Step Functions (Standard)
    & Слоевете Planning, Execution, Aggregation, Classification и Operation;
      паралелно изпълнение на агентите; история на всяко разследване за
      одит\\ \hline
  Изчисления & Lambda (в VPC)
    & Агентите, обединяването, класификацията и интеграциите\\ \hline
  Изчисления & ECS Fargate
    & Playwright сценариите и n8n (зад вътрешен ALB)\\ \hline
  LLM & Amazon Bedrock
    & Езиковите модели, използвани от ЕС регион; данните не излизат от
      AWS\\ \hline
  Горещи данни & ElastiCache Serverless (Valkey)
    & Общ контекст на агентите, прозорец за дедупликация, ограничаване на
      заявките към външни API\\ \hline
  Състояние & DynamoDB
    & Инциденти, алърти и статус на разследванията; данни за таблото\\ \hline
  История & Aurora PostgreSQL Serverless v2 с pgvector
    & История на инцидентите, търсене на сходни случаи, мапинг на
      компонентите\\ \hline
  Архив & S3
    & Файловете на таблото, сурови събития, отчети и данни за обучение
      на моделите\\ \hline
  Тайни & Secrets Manager и KMS
    & Токените за Splunk, тикет системата и GCP; криптиране на данните\\ \hline
  Наблюдение & CloudWatch, X-Ray, OpenTelemetry
    & Логове, метрики и трасировки на самата YTKI; аларми\\ \hline
\end{xltabular}
\endgroup

\subsection*{А.6. Мрежа и сигурност}

\textbf{Регион.} Цялата система работи в eu-central-1 (Франкфурт). Данните
остават в ЕС, което улеснява спазването на GDPR.

\textbf{VPC.} Използваме три зони на достъпност. Lambda функциите, задачите
във Fargate, ElastiCache и Aurora са в частни подмрежи без публични
адреси. Публичните подмрежи съдържат само NAT Gateway. Lambda функциите
се свързват с VPC чрез мрежови интерфейси (ENI), които AWS създава при
конфигуриране на функцията, затова подмрежите трябва да имат достатъчно
свободни адреси. Security групите разрешават достъп до Aurora и ElastiCache
само от функциите и задачите на YTKI.

\textbf{Изходящ трафик.} Към S3 и DynamoDB се стига през безплатни gateway
endpoints, а към Bedrock, SQS, Step Functions, Secrets Manager, KMS,
CloudWatch и ECR през interface endpoints. Така трафикът към AWS услугите
не минава през NAT. NAT Gateway остава само за външни системи като GCP,
Splunk и тикет системата. Ако тези системи поддържат IPv6, NAT може да се
замени с egress-only internet gateway, който няма почасова такса.

\textbf{Входна точка за алъртите.} Всеки източник получава отделен токен,
който се проверява от Lambda authorizer. Към API-то е закачен WAF
с ограничение на броя заявки. Като следваща стъпка може да се добави mTLS.

\textbf{Достъп до GCP.} Агентът за логове чете от Cloud Logging чрез
Workload Identity Federation: IAM ролята на функцията се разменя за
краткотраен токен в GCP, без да пазим ключове.

\textbf{Права.} Всяка функция има собствена IAM роля само с най-необходимите
права. CI/CD конвейерът влиза в AWS чрез OIDC, без постоянни ключове.
CloudTrail, GuardDuty и Security Hub са включени във всички акаунти.

\textbf{Данни.} Всичко се криптира с ключове в KMS. Личните данни се
маскират преди подаване към Bedrock (задача 2.6), а логовете на самата
YTKI се пазят ограничено време.

\subsection*{А.7. Среди и стратегии за внедряване}

\textbf{Акаунти.} Чрез AWS Organizations имаме отделни акаунти за
инструментите (CI/CD, ECR), за dev, staging и prod на YTKI и за
централизиран архив на логовете. Средите на YTKI не съвпадат със средите,
които тя наблюдава: prod версията на YTKI приема алърти от всички четири
наблюдавани среди, а dev и staging версиите получават копия или записани
исторически алърти.

\textbf{Инфраструктура като код.} Цялата инфраструктура е описана
с Terraform в модули (мрежа, данни, изчисления, край на мрежата). Едни
и същи модули се използват за всички акаунти, с различни параметри.
Избираме Terraform, защото с него може да опишем и частта в GCP
(Workload Identity Federation), а не само тази в AWS.

\textbf{Конвейер.} GitHub Actions с OIDC към AWS. Стъпките са: проверки
и unit тестове, изграждане на пакетите за Lambda, образите за ECR
и файловете на таблото, \texttt{terraform plan} с ревю в pull request,
внедряване в dev, интеграционни тестове с повторно изпращане на исторически
алърти, внедряване в staging, ръчно одобрение и внедряване в prod.

\begingroup\small
\begin{xltabular}{\textwidth}{|P{3.1cm}|L|P{4.2cm}|}
  \caption{Стратегии за пускане на нови версии}\label{tab:deploy}\\
  \hline
  \rowcolor{hdr}
  \textbf{Компонент} & \textbf{Стратегия} & \textbf{Връщане назад}\\ \hline
  \endfirsthead
  \hline
  \rowcolor{hdr}
  \textbf{Компонент} & \textbf{Стратегия} & \textbf{Връщане назад}\\ \hline
  \endhead
  Lambda функции
    & Версии и псевдоними; CodeDeploy пренасочва първо 10\% от трафика
      и следи алармите
    & Автоматично, ако аларма се задейства\\ \hline
  Step Functions
    & Версии и псевдоними с постепенно пренасочване между старата и новата
      версия
    & Псевдонимът се връща към старата версия\\ \hline
  Fargate (n8n, Playwright)
    & Вградения blue/green в ECS
    & Трафикът се връща към старите задачи\\ \hline
  Табло (S3 и CloudFront)
    & Новата версия се качва в отделна папка и CloudFront се пренасочва след
      проверка
    & Пренасочване към предишната папка\\ \hline
  Схема на Aurora
    & Миграции в две стъпки (expand/contract): първо се добавя новото,
      старото се маха в следваща версия
    & Старият код продължава да работи с новата схема\\ \hline
  Класификатор и промптове
    & Сенчест режим: новата версия работи паралелно със старата, без да
      предприема действия; резултатите се сравняват
    & Новата версия просто не се активира\\ \hline
  Инфраструктура
    & \texttt{terraform plan} в pull request, ревю, прилагане само през
      конвейера
    & Прилагане на предишната версия на кода\\ \hline
\end{xltabular}
\endgroup

\subsection*{А.8. Надеждност и наблюдение на самата система}

\textbf{YTKI не е в критичния път.} Алъртите стигат до дежурните по
стария път, а YTKI само ги допълва с анализ. Ако тя спре, губим анализа,
но не и алъртите.

\textbf{Нищо не се губи.} SQS опашката с DLQ и архивът в EventBridge
позволяват алъртите да се изпратят повторно. Отпечатъкът на алърта служи
като ключ за идемпотентност, затова повторното изпращане не създава
дублирани тикети.

\textbf{Контролен сигнал.} На всеки пет минути през системата минава
изкуствен алърт. Ако за него не се появи отчет навреме, CloudWatch вдига
аларма директно към PagerDuty.

\textbf{Отказ на LLM.} При грешки или ограничения от Bedrock Step Functions
повтаря заявката с нарастващо изчакване, а при продължителен проблем
алъртът минава без анализ, с ясна бележка в таблото.

\textbf{Резервни копия.} DynamoDB е с възстановяване към точен момент,
Aurora прави автоматични резервни копия, а S3 пази версии. AWS Backup
копира данните и в друг ЕС регион. Предлагаме цел за загуба на данни до
15 минути (RPO) и за възстановяване до 4 часа (RTO).

\textbf{Наблюдение.} Всяко разследване се проследява с OpenTelemetry от
приемането до действието. Таблата в CloudWatch показват време за
разследване, грешки, изразходвани токени и цена на един инцидент.

\subsection*{А.9. Разходи}

Променливите разходи са за Lambda, Step Functions, API Gateway и най-вече
за токените в Bedrock. Постоянните разходи са за NAT Gateway и interface
endpoints (таксуват се на час за всяка зона), минималния капацитет на
Aurora и ElastiCache и постоянно работещата задача на n8n. За да ги
намалим, в dev средата използваме една зона и един NAT, а Aurora спира,
когато не се използва. Дубликатите се отсяват преди LLM, за по-простите
стъпки (планиране, класификация) се ползва по-малък модел, а AWS Budgets
и Cost Anomaly Detection следят сметката. Точната оценка ще бъде част от
т.~2.3 (разпределение на ресурсите).

\subsection*{А.10. Разгледани алтернативи}

\textbf{Lambda durable functions вместо Step Functions.} Позволяват
многостъпкови процеси директно в кода, без отделна услуга. Предпочитаме
Step Functions, защото дава визуална история на всяко разследване, която
е полезна за одит, и истинско паралелно изпълнение на агентите.

\textbf{Amazon Bedrock AgentCore Runtime за агентите.} Предлага изолирани
сесии с продължителност до няколко часа. Подходящ е, ако „Ask Alerts“
прерасне в дълги интерактивни разследвания, и е възможна следваща стъпка
след v1.0.

\textbf{Само DynamoDB, без ElastiCache.} Ще намали броя на компонентите.
Запазваме ElastiCache заради бързия общ контекст и ограничаването на
заявките при алърт бури, но в първата версия тя може да се пропусне
и да се добави, ако измерванията покажат нужда.

\end{document}
